\begin{abstract}
O Sistema Integrado de Gestão de Atividades Acadêmicas (SIGAA) é o ponto único de acesso do estudante da Universidade de Brasília (UnB) a notas, horários, histórico, documentos e turmas virtuais. Relatos de dificuldade de uso, uma interface pensada para computador de mesa e ferramentas comunitárias que suprem funções do sistema motivam uma avaliação objetiva da sua qualidade. Este artigo apresenta a especificação dessa avaliação, segundo o processo de avaliação da família SQuaRE (ISO/IEC 25040) e o modelo de qualidade da ISO/IEC 25010:2023. Na Fase 1, uma matriz Impacto~$\times$~Risco com filtro de viabilidade priorizou quatro características, e personas deram concretude aos perfis de usuários. Por orientação recebida no ponto de controle, a especificação concentra-se em três delas: Capacidade de Interação (Usabilidade), Adequação Funcional e Flexibilidade (Portabilidade). Na Fase 2, o método \textit{Goal-Question-Metric} (GQM) as desdobrou em três objetivos, 12 questões com hipóteses e 16 métricas, com níveis de pontuação e critérios de julgamento fixados antes da coleta. A especificação combina teste com usuários, inspeção por dois avaliadores, verificação automática de acessibilidade e um questionário que compara estudantes de diferentes semestres. A gestão do projeto é acompanhada por indicadores do PSM/CID.
\end{abstract}

\begin{IEEEkeywords}
qualidade de produto de software, SQuaRE, ISO/IEC 25010, GQM, PSM, SIGAA.
\end{IEEEkeywords}
